\section{Convergence, remainder estimates, and logarithm obstructions}\label{sec:convergence}
The formal theorem is unconditional, but its numerical evaluation is not.
Four different questions must be kept apart: convergence of the Mercator
series $\sum_k(-1)^{k-1}(e^Xe^Y-I)^k/k$; absolute convergence of its
expansion into individual associative words or into the individual
commutator summands of Dynkin's formula; convergence of the homogeneous Lie
series $\sum_nZ_n$; and existence of some logarithm of $e^Xe^Y$. These are
different summation procedures, and outside a common domain of absolute
convergence they can behave differently; examples and further discussion
appear in \cite{biagi}. This section proves elementary quantitative
statements sufficient for every convergence claim on the source page,
without claiming optimal constants.

\subsection{Absolute convergence in a Banach algebra}
Let $\cA$ be a real or complex unital Banach algebra with a submultiplicative
norm, $\norm{AB}\leq\norm A\norm B$. Unless said otherwise we do not assume
$\norm I=1$: every estimate below involves only products of positive-length
words, so the norm of the identity never enters. In particular the
Hilbert--Schmidt norm on $d\times d$ matrices is admissible: by
Cauchy--Schwarz,
\[
 \norm{AB}_{\mathrm{HS}}^2=\sum_{i,j}\Bigl|\sum_kA_{ik}B_{kj}\Bigr|^2
 \leq\sum_{i,j}\Bigl(\sum_k|A_{ik}|^2\Bigr)\Bigl(\sum_k|B_{kj}|^2\Bigr)
 =\norm A_{\mathrm{HS}}^2\norm B_{\mathrm{HS}}^2,
\]
although $\norm I_{\mathrm{HS}}=\sqrt d$. Throughout, $s=\norm X+\norm Y$.

\begin{theorem}[Two convergence domains and explicit remainders]\label{thm:convergence}
\begin{enumerate}[label=(\roman*)]
\item If $s<\log2$, the associative block expansion \eqref{eq:blocks} is
absolutely convergent, the homogeneous BCH series $\sum_nZ_n(X,Y)$ converges
absolutely, and its sum $Z$ satisfies $e^Z=e^Xe^Y$. More precisely, for
every $r>0$ with $rs<\log2$,
\begin{equation}\label{eq:majorant}
\begin{gathered}
 \sum_{k\geq1}\frac1k\sum_{\substack{r_i+s_i>0}}
 \frac{\norm{X^{r_1}Y^{s_1}\cdots X^{r_k}Y^{s_k}}}{\prod_ir_i!s_i!}\,r^{\sum(r_i+s_i)}
 \leq M_s(r):=-\log\bigl(2-e^{rs}\bigr),\\
 \sum_{n\geq1}r^n\norm{Z_n(X,Y)}\leq M_s(r).
\end{gathered}
\end{equation}
\item If $1<r<\log2/s$ and $Z^{[N]}=\sum_{n=1}^NZ_n$, then
\begin{align}
 \norm{Z-Z^{[N]}}&\leq r^{-(N+1)}M_s(r),\label{eq:remainder}\\
 \norm{e^{Z^{[N]}}-e^Xe^Y}&\leq c_{\cA}\,e^{M_s(1)}\,r^{-(N+1)}M_s(r),
 \label{eq:expremainder}
\end{align}
where $c_{\cA}=\norm I^2$; thus $c_{\cA}=1$ for a norm with $\norm I=1$.
\item If $s<\tfrac12\log2$, the individual right-nested commutator summands
of Dynkin's formula \eqref{eq:dynkin} are absolutely summable, with
\begin{equation}\label{eq:dynkinradius}
 \sum_{\text{all Dynkin summands}}\norm{\text{summand}}
 \leq M_D(s):=-\tfrac12\log\bigl(2-e^{2s}\bigr),
\end{equation}
and if $Rs<\tfrac12\log2$ with $R>1$ the tail after degree $N$ of this
absolute sum is at most $R^{-(N+1)}M_D(Rs)$.
\item For a finite-dimensional normed Lie algebra $\g$ with
$\norm{[u,v]}\leq\kappa\norm u\norm v$, Dynkin's Lie series converges
absolutely whenever $\kappa(\norm X+\norm Y)<\log2$, with absolute sum at
most $\kappa^{-1}\bigl[-\log(2-e^{\kappa s})\bigr]$.
\end{enumerate}
\end{theorem}
\begin{proof}
(i) Expanding into nonempty words,
\[
 \norm{e^Xe^Y-I}\leq\sum_{r+s>0}\frac{\norm X^r\norm Y^s}{r!s!}=e^{s}-1,
\]
where a block with $r=0$ or $s=0$ is a pure power and contributes no factor
$\norm I$. Replacing each term of the block expansion of $(e^Xe^Y-I)^k$ by
the product of the norms of its blocks bounds the sum of the norms of all
terms of the $k$th power by $(e^s-1)^k$. If $s<\log2$ then $e^s-1<1$, and
\[
 \sum_{k\geq1}\frac{(e^s-1)^k}{k}=-\log\bigl(1-(e^s-1)\bigr)=-\log(2-e^s).
\]
Replacing $X,Y$ by $rX,rY$ multiplies a term of total degree $n$ by $r^n$
and gives the first inequality in \eqref{eq:majorant}. Absolute convergence
of the expanded series permits regrouping by total degree; the group of
degree $n$ is $Z_n(X,Y)$ by \cref{thm:wordcut} (the formal identity is a
statement about which word coefficients occur), so
$\sum_nr^n\norm{Z_n}\leq M_s(r)$, and with $r=1$ the homogeneous series
converges absolutely. Its sum is the Mercator logarithm $\log(I+U)$,
$U=e^Xe^Y-I$, $\norm U<1$. The scalar identity $\exp(\log(1+u))=1+u$ holds
for $|u|<1$ as an identity of absolutely convergent series, and substituting
$u\mapsto U$ is legitimate in the closed commutative subalgebra generated by
$U$, exactly as in \cref{lem:explog}. Hence $e^Z=I+U=e^Xe^Y$.

(ii) For $n>N$, $r^n\geq r^{N+1}$, so
$\sum_{n>N}\norm{Z_n}\leq r^{-(N+1)}\sum_{n>N}r^n\norm{Z_n}\leq r^{-(N+1)}M_s(r)$,
which is \eqref{eq:remainder}. For the exponentials, differentiating
$u\mapsto e^{(1-u)A}e^{uB}$ gives
\begin{equation}\label{eq:expdifference}
 e^B-e^A=\int_0^1e^{(1-u)A}(B-A)e^{uB}\,\dd u,
\end{equation}
since the derivative of the integrand's antiderivative
$e^{(1-u)A}e^{uB}$ is $e^{(1-u)A}(-A+B)e^{uB}$ (the factor $A$ commutes with
$e^{(1-u)A}$ and $B$ with $e^{uB}$). Submultiplicativity gives $\norm I=\norm{I^2}\leq\norm I^2$, so
$\norm I\geq1$, and therefore
$\norm{e^C}\leq\norm I+\sum_{n\geq1}\norm C^n/n!\leq\norm I\,e^{\norm C}$.
Consequently
\[
 \norm{e^B-e^A}\leq\norm I^2\int_0^1e^{(1-u)\norm A}e^{u\norm B}\dd u\,\norm{B-A}
 \leq\norm I^2e^{\max(\norm A,\norm B)}\norm{B-A}.
\]
Apply this with $A=Z^{[N]}$ and $B=Z$: both have norm at most
$\sum_n\norm{Z_n}\leq M_s(1)$, and $\norm{Z-Z^{[N]}}$ is bounded by
\eqref{eq:remainder}. Since $e^Z=e^Xe^Y$ by (i), this proves
\eqref{eq:expremainder}.

(iii) For a right-nested bracket of length $n$, iterating
$\norm{[U,V]}\leq2\norm U\norm V$ gives
$\norm{R(a_1\cdots a_n)}\leq2^{n-1}\prod_j\norm{a_j}$. In \eqref{eq:dynkin}
discard the factor $1/d\leq1$. The sum of the norms of the degree-$n$
summands is then at most $2^{n-1}$ times the sum of the norms of the
corresponding associative block terms with $X,Y$ replaced by $2X,2Y$, divided
by $2^n$; hence the total is at most
$\frac12\sum_{k\geq1}(e^{2s}-1)^k/k=-\frac12\log(2-e^{2s})$, finite when
$2s<\log2$. The tail bound follows exactly as in (ii) by applying the
estimate to $RX,RY$.

(iv) Replace $2^{n-1}$ by $\kappa^{n-1}$ in (iii); the norm bound on a
right-nested bracket of length $n$ becomes $\kappa^{n-1}\prod_j\norm{a_j}$,
and the same computation gives the bound $\kappa^{-1}[-\log(2-e^{\kappa s})]$
under $\kappa s<\log2$. Every finite-dimensional Lie algebra admits such a
$\kappa$, because a bilinear map is bounded on the product of the unit
spheres, so this proves local convergence without a matrix representation.
If the bracket is identically zero, BCH terminates at $X+Y$ and no
positive $\kappa$ is needed.
\end{proof}

The two radii $\log2$ and $\tfrac12\log2$ concern different ways of summing
expressions. The larger applies to the series of already-collected
homogeneous Lie polynomials and to the associative word expansion; the
smaller controls all the unreduced nested-commutator terms of Dynkin's
multiple sum separately, and it is the bound stated in the source page's
convergence discussion. Neither is claimed to be optimal.

\subsection{The Mercator tail and the role of \texorpdfstring{$\norm Z<\log2$}{||Z||<log 2}}
For a direct truncation of the Mercator series, put $q=e^s-1<1$. Bounding
$1/k\leq1/(K+1)$ in the tail gives
\begin{equation}\label{eq:mercatortail}
 \Bigl\|\log(e^Xe^Y)-\sum_{k=1}^K\frac{(-1)^{k-1}}{k}(e^Xe^Y-I)^k\Bigr\|
 \leq\sum_{k>K}\frac{q^k}{k}\leq\frac{q^{K+1}}{(K+1)(1-q)}.
\end{equation}
The source page displays the additional condition $\norm Z<\log2$ next to
its associative expansion. That condition is \emph{not} needed to define
$Z=\log(e^Xe^Y)$ by the convergent Mercator series: the bound actually
proved is $\norm Z\leq M_s(1)$, and one should not silently infer a
further norm bound on this constructed $Z$. The condition is relevant in
the converse direction. Suppose $Z_0$ is \emph{given} with $e^{Z_0}=e^Xe^Y$
and $\norm{Z_0}<\log2$. Then for $|t|\norm{Z_0}<\log2$ we have
$\norm{e^{tZ_0}-I}\leq e^{|t|\norm{Z_0}}-1<1$, so the Mercator logarithm
$L(t)=\log(e^{tZ_0})$ is a convergent power series in $t$ on that disk. Near
$t=0$ it equals $tZ_0$ by the formal identity $\log(e^{u})=u$ applied in the
commutative closed subalgebra generated by $Z_0$; both sides being analytic
in $t$ on the disk, they agree throughout it, and at $t=1$ we get
$\log(e^{Z_0})=Z_0$. Combined with $s<\log2$, this shows that the page's two
restrictions do suffice to conclude that its given $Z_0$ coincides with the
Mercator logarithm of $e^Xe^Y$, and hence with the BCH sum. Without a
restriction on $Z_0$ the conclusion fails: other logarithms exist, as
the trace discussion next shows.

\subsection{The trace identity and its branch restriction}
\begin{proposition}\label{prop:trace}
For the formal BCH series, and for finite complex or real matrices whenever
the homogeneous series converges,
\begin{equation}\label{eq:trace}
 \tr Z(X,Y)=\tr X+\tr Y.
\end{equation}
For an arbitrary complex matrix $L$ with $e^L=e^Xe^Y$, the branch-independent
conclusion is only
\begin{equation}\label{eq:tracebranch}
 \tr L-\tr X-\tr Y\in2\pi\ii\Z.
\end{equation}
\end{proposition}
\begin{proof}
The identity $\tr(AB)=\tr(BA)$ follows by interchanging the two summation
indices in $\sum_{i,j}A_{ij}B_{ji}$. Hence every commutator has trace zero,
and every nested bracket of length at least two, being a commutator, has
trace zero. By \cref{thm:formalbch} each $Z_n$ with $n\geq2$ is a sum of
such brackets, so $\tr Z_n=0$ for $n\geq2$ and $\tr Z_1=\tr X+\tr Y$. This
proves \eqref{eq:trace} degree by degree; when the series converges,
continuity of the trace (a linear functional on a finite-dimensional space)
gives it for the sum.

For \eqref{eq:tracebranch} we first show $\det e^A=e^{\tr A}$ for every
complex square matrix $A$. Every complex matrix is triangularizable: choose
an eigenvector $v_1$, complete it to a basis, and observe that in this basis
$A$ is block upper triangular with a $1\times1$ block; induct on the
dimension of the quotient. In an upper-triangular basis with diagonal
$\lambda_1,\ldots,\lambda_d$, every power $A^k$ is upper triangular with
diagonal $\lambda_j^k$, so $e^A$ is upper triangular with diagonal
$e^{\lambda_j}$, and $\det e^A=\prod_je^{\lambda_j}=e^{\sum\lambda_j}
=e^{\tr A}$. Applying this to $e^L=e^Xe^Y$ and using multiplicativity of
the determinant gives $e^{\tr L}=e^{\tr X}e^{\tr Y}$, that is
$e^{\tr L-\tr X-\tr Y}=1$, and the scalar exponential takes the value $1$
exactly on $2\pi\ii\Z$.
\end{proof}
A scalar counterexample to unrestricted use of the principal logarithm:
take $X=Y=3\pi\ii/4$ in dimension one. The BCH logarithm is $X+Y=3\pi\ii/2$,
whereas the principal logarithm of the product $e^Xe^Y=e^{3\pi\ii/2}=-\ii$
is $-\pi\ii/2$; they differ by $2\pi\ii$. Thus the word ``logarithm'' in
\eqref{eq:trace} cannot be left branch-independent. When $L,X,Y$ are all
real matrices the difference in \eqref{eq:tracebranch} is real, hence zero.

\subsection{The \texorpdfstring{$\mathfrak{sl}_2$}{sl2} example, completely resolved}
\begin{theorem}\label{thm:sl2}
Let
\begin{equation}\label{eq:sl2XY}
 X=\begin{pmatrix}0&\ii\pi\\\ii\pi&0\end{pmatrix},\qquad
 Y=\begin{pmatrix}0&1\\0&0\end{pmatrix}.
\end{equation}
Then
\begin{equation}\label{eq:sl2product}
 e^Xe^Y=\begin{pmatrix}-1&-1\\0&-1\end{pmatrix}=-I-Y.
\end{equation}
The complete set of complex $2\times2$ logarithms of this matrix is
\begin{equation}\label{eq:sl2logs}
 L_k=(2k+1)\pi\ii\,I+Y,\qquad k\in\Z.
\end{equation}
None of them is traceless, so none lies in $\mathfrak{sl}_2(\C)$, and the
homogeneous BCH series $\sum_nZ_n(X,Y)$ does not converge, even
conditionally, although the product has infinitely many logarithms in
$\mathfrak{gl}_2(\C)$.
\end{theorem}
\begin{proof}
Since $X^2=-\pi^2I$, the even and odd parts of the exponential series give
$e^X=\sum_m\frac{(-\pi^2)^m}{(2m)!}I+\sum_m\frac{(-\pi^2)^m}{(2m+1)!}X
=\cos\pi\,I+\frac{\sin\pi}{\pi}X=-I$. Since $Y^2=0$, $e^Y=I+Y$. Hence
$e^Xe^Y=-I-Y$, which is \eqref{eq:sl2product}.

Let $L$ be any logarithm, $e^L=-I-Y$. A matrix commutes with every power of
itself, hence with its exponential; so $L$ commutes with $-I-Y$ and
therefore with $Y$. Writing $L=\bigl(\begin{smallmatrix}a&b\\c&d\end{smallmatrix}\bigr)$,
the equation $LY=YL$ reads
$\bigl(\begin{smallmatrix}0&a\\0&c\end{smallmatrix}\bigr)
=\bigl(\begin{smallmatrix}c&d\\0&0\end{smallmatrix}\bigr)$, so $c=0$ and
$a=d$: the centralizer of $Y$ consists of the matrices $aI+bY$. For such a
matrix, $aI$ and $bY$ commute and $(bY)^2=0$, so
$e^{aI+bY}=e^a(I+bY)$. Equality with $-I-Y$ forces $e^a=-1$ and $e^ab=-1$,
that is $a=(2k+1)\pi\ii$ and $b=1$. This proves \eqref{eq:sl2logs}. The
trace of $L_k$ is $2(2k+1)\pi\ii\neq0$.

Suppose, to reach a contradiction, that $\sum_{n\geq1}Z_n(X,Y)$ converged to
some matrix $S$. Then $Z_n\to0$, so the coefficients are bounded and the
matrix power series $Z(t)=\sum_{n\geq1}t^nZ_n(X,Y)$ converges absolutely,
entry by entry, for $|t|<1$, defining an analytic matrix function there. For
$|t|(\norm X+\norm Y)<\log2$ we have $Z(t)=\BCH(tX,tY)$ by homogeneity, and
\cref{thm:convergence}(i) gives $e^{Z(t)}=e^{tX}e^{tY}$. Both sides are
analytic in $t$ on the unit disk (the exponential of an analytic matrix
function is analytic, being an absolutely and locally uniformly convergent
series of analytic functions), so by the identity theorem
$e^{Z(t)}=e^{tX}e^{tY}$ for all $|t|<1$.

Abel's limit theorem now gives $Z(t)\to S$ as $t\uparrow1$ along the real
axis. For completeness: let $S_m=\sum_{n=1}^mZ_n$ and $S_0=0$. For $0<t<1$,
summation by parts gives
\[
 Z(t)=\sum_{n\geq1}t^n(S_n-S_{n-1})=\sum_{n\geq1}t^nS_n-\sum_{n\geq1}t^{n+1}S_n
 =(1-t)\sum_{n\geq1}t^nS_n,
\]
all series converging absolutely because the $S_n$ are bounded. Since
$(1-t)\sum_{n\geq1}t^n=t$, we obtain
$Z(t)-tS=(1-t)\sum_{n\geq1}t^n(S_n-S)$. Given $\eps>0$ choose $N$ with
$\norm{S_n-S}<\eps$ for all $n>N$. The finite part
$(1-t)\sum_{n\leq N}t^n(S_n-S)$ tends to zero as $t\uparrow1$, and the tail
is bounded by $\eps(1-t)\sum_{n>N}t^n\leq\eps$. Hence
$\limsup_{t\uparrow1}\norm{Z(t)-tS}\leq\eps$ for every $\eps>0$, so
$Z(t)-tS\to0$ and $Z(t)\to S$. Continuity of the exponential then gives $e^S=\lim e^{Z(t)}
=\lim e^{tX}e^{tY}=e^Xe^Y=-I-Y$. But every $Z_n$ is traceless by
\cref{prop:trace}, so $\tr S=\lim\tr S_m=0$, contradicting the
classification \eqref{eq:sl2logs}. Hence the series does not converge.
\end{proof}
This example simultaneously distinguishes existence of a logarithm,
existence of a logarithm in the specified Lie algebra, and convergence of
the BCH series. They are genuinely different questions. The obstruction is
not that the invertible matrix $-I-Y$ lacks a logarithm; it has infinitely
many. The obstruction is that it has no logarithm in the traceless Lie
algebra generated by $X$ and $Y$.
